\begin{frame}[t]{Forkable sandboxes let computers explore executable alternatives.}
\vspace{0pt}

\vspace{.15cm}
{\fontsize{21}{25}\selectfont\color{Teal}Self-organization $\to$ self-driving cars\par $\to$ self-driving compute.\par}
\vspace{.35cm}
\forktree\quad {\fontsize{14}{17}\selectfont Reuse reached state. Run alternatives. Verify the result.\par}
\vspace{.35cm}
{\fontsize{12.5}{15}\selectfont Yossi Eliaz\par Principal Engineer, Incredibuild\par Associate Professor, HIT\par}

\sources{Cambridge Systems Research Group, 15 October 2026. Research and engineering perspectives; employer affiliation disclosed.}
\qadetail{The integrated architecture is a research design. Individual implementations, published examples and own measurements are labeled on their respective slides. Employer affiliation does not imply endorsement. The title figure is a proposed lifecycle. A motivating future workload is a running application at a reproduced failure, with relevant process memory and local state. It has not been exercised in the factory case study. The observed serial repair establishes source files and feedback, not a demonstrated need for live fork. Select a worktree, warm image, disk snapshot, or verified process/VM checkpoint according to the required state. Disclosure and affiliations match the supplied Cambridge deck.}
\note{[0:25] Once we have paid to reach a useful computational state, can we explore several next steps from there? I will connect my earlier work to this question, explain the execution model through examples, and examine what each continuation must return for its result to be accepted.}
\end{frame}
